% !TEX program = lualatex
\documentclass[11pt,a4paper]{article}
\usepackage{luatexja}
\usepackage{amsmath,amssymb,amsthm}
\usepackage{geometry}
\geometry{margin=1in}
\usepackage{booktabs}
\usepackage{array}
\usepackage{colortbl}
\usepackage[table]{xcolor}

%-------------------------------------------------
\newcommand{\mweq}{\mathrel{\overset{\mathrm{mw}}{=}}}
\newcommand{\dfix}{D_{\mathrm{fix}0}}

\theoremstyle{definition}
\newtheorem{definition}{定義}[section]
\newtheorem{theorem}{定理}[section]
\newtheorem{remark}{備考}[section]

\renewcommand{\abstractname}{Abstract}

%-------------------------------------------------
\begin{document}

\title{再帰的四句分別とグレイコード：\\16真理関数との対応表}
\author{千葉将臣}
\date{2026-04-28}
\maketitle

\begin{abstract}
四句分別の再帰的展開（$4^2 = 16$）と古典論理の16真理関数は、
同じ濃度を持ちながら異なる順序原理に基づいて並んでいる。
本稿は、四句分別がグレイコード的変化（隣接項の1ビット変化）を持つのに対し、
16真理関数が順序数的変化（バイナリの逐次増加）を持つことを示し、
両者の対応表をグレイコード／バイナリ変換として導出する。
\end{abstract}

\section{順序原理の差異}

\begin{definition}[グレイコード的変化]
系列における隣接項が、ビット表現において常に1ビットのみ異なる変化様式。
境界の最小変化に対応する。
\end{definition}

\begin{definition}[順序数的変化]
$0, 1, 2, \ldots, 2^n - 1$ というバイナリの逐次増加による変化様式。
自己同一性を前提とした線形順序に対応する。
\end{definition}

四句分別の系列における変化は以下の通りである：

\[
\text{是}(00) \;\to\; \text{非}(01) \;\to\; \text{亦是亦非}(11) \;\to\; \text{非是非非}(10)
\]

隣接項の差分はそれぞれ1ビットであり、グレイコードの定義を満たす。

一方、16真理関数は入力パターン $(T,T),(T,F),(F,T),(F,F)$ への出力割り当てとして
$f_0, f_1, \ldots, f_{15}$（$0 = 0000_2$ から $15 = 1111_2$）という
順序数的列挙として導入される。

\section{グレイコード／バイナリ変換による対応}

四句値をグレイコードに対応させる：

\begin{center}
\begin{tabular}{ccc}
\toprule
四句値 & グレイコード & バイナリ（順序数） \\
\midrule
是      & \texttt{00} & 0 \\
非      & \texttt{01} & 1 \\
亦是亦非 & \texttt{11} & 2 \\
非是非非 & \texttt{10} & 3 \\
\bottomrule
\end{tabular}
\end{center}

外軸・内軸のグレイコードを連結した4ビット値をバイナリに変換することで、
対応する真理関数の順序数が得られる。

\section{対応表}

\renewcommand{\arraystretch}{1.25}
\begin{center}
\begin{tabular}{llcccp{4.5cm}}
\toprule
外軸 & 内軸 & 連結GC & 順序数 & 真理関数 & 意味 \\
\midrule
\rowcolor{gray!10}
是 & 是     & \texttt{0000} & 0  & $f_0$  & $\bot$（恒偽） \\
是 & 非     & \texttt{0001} & 1  & $f_1$  & NOR $(\downarrow)$：$\neg P \wedge \neg Q$ \\
\rowcolor{gray!10}
是 & 亦是亦非 & \texttt{0011} & 2  & $f_2$  & $\neg P \wedge Q$ \\
是 & 非是非非 & \texttt{0010} & 3  & $f_3$  & $\neg P$ \\
\midrule
\rowcolor{gray!10}
非 & 是     & \texttt{0100} & 7  & $f_7$  & NAND $(\uparrow)$：$\neg(P \wedge Q)$ \\
非 & 非     & \texttt{0101} & 6  & $f_6$  & XOR $(\oplus)$：$P \oplus Q$ \\
\rowcolor{gray!10}
非 & 亦是亦非 & \texttt{0111} & 5  & $f_5$  & $\neg Q$ \\
非 & 非是非非 & \texttt{0110} & 4  & $f_4$  & $P \wedge \neg Q$ \\
\midrule
\rowcolor{gray!10}
亦是亦非 & 是     & \texttt{1100} & 8  & $f_8$  & AND $(\wedge)$ \\
亦是亦非 & 非     & \texttt{1101} & 9  & $f_9$  & XNOR $(\leftrightarrow)$：$P \leftrightarrow Q$ \\
\rowcolor{gray!10}
亦是亦非 & 亦是亦非 & \texttt{1111} & 10 & $f_{10}$ & $Q$（射影） \\
亦是亦非 & 非是非非 & \texttt{1110} & 11 & $f_{11}$ & $P \to Q$（含意） \\
\midrule
\rowcolor{gray!10}
非是非非 & 是     & \texttt{1000} & 15 & $f_{15}$ & $\top$（恒真） \\
非是非非 & 非     & \texttt{1001} & 14 & $f_{14}$ & OR $(\vee)$ \\
\rowcolor{gray!10}
非是非非 & 亦是亦非 & \texttt{1011} & 13 & $f_{13}$ & $Q \to P$（逆含意） \\
非是非非 & 非是非非 & \texttt{1010} & 12 & $f_{12}$ & $P$（射影） \\
\bottomrule
\end{tabular}
\end{center}

\section{対応表の構造的観察}

\begin{theorem}[順序数の非単調性]
四句分別マトリクスの行順に読んだ真理関数の順序数列は
$0,1,2,3,7,6,5,4,8,9,10,11,15,14,13,12$
であり、単調増加しない。
各ブロック内では単調増加または減少が交互に現れる。
\end{theorem}

\begin{remark}[非単調性の意味]
真理関数の順序数的並びは $id$ を前提とした線形順序であり、
四句分別マトリクスのグレイコード的並びはフラクタル境界の最小変化に対応する。
両者の対応が非単調になることは、二つの順序原理が根本的に異なることの
直接的な反映である。
\end{remark}

\begin{remark}[$\bot$ と $\top$ の位置]
順序数的並びでは $\bot = f_0$（先頭）、$\top = f_{15}$（末尾）であり、
対称的な両端に位置する。
四句分別マトリクスでは $\bot$（是×是）と $\top$（非是非非×是）が
対角的な位置に現れ、線形な対称性を持たない。
これは $\top, \bot$ が $id$ を前提とした線形順序の両端ではなく、
フラクタル構造の特定の位置として現れることを示している。
\end{remark}

\section{変換の界論的意味}

グレイコード／バイナリ変換は、四句分別マトリクスから真理関数への
座標変換として機能する。

\begin{itemize}
  \item グレイコード座標：境界の最小変化を基準とするフラクタル的記述
  \item バイナリ座標：$id$ を前提とした線形順序による記述
\end{itemize}

二元論（古典論理）が真理関数を順序数として列挙したのは、
グレイコード座標の存在を知らずにバイナリ座標のみを採用したからである。
界論的には、グレイコード座標が生成機序に対応し、
バイナリ座標はその静的断面である。

\end{document}
